\documentclass[11pt,twocolumn]{article}

% -----------------------------
% PACKAGES (SAFE FOR OVERLEAF)
% -----------------------------
\usepackage[margin=1in]{geometry}
\usepackage[T1]{fontenc}
\usepackage[utf8]{inputenc}
\usepackage{lmodern}
\usepackage{amsmath,amssymb}
\usepackage{graphicx}
\usepackage{booktabs}
\usepackage{enumitem}
\usepackage{hyperref}
\usepackage{float}
\usepackage{caption}
\usepackage{titlesec}
\usepackage{url}
\usepackage{array}

% -----------------------------
% SPACING OPTIMIZATION
% -----------------------------
\setlength{\columnsep}{0.25in}
\setlength{\parindent}{1em}
\setlength{\parskip}{0pt}
\setlist[itemize]{itemsep=0pt, topsep=0pt, parsep=0pt, partopsep=0pt}

\titlespacing*{\section}{0pt}{0.7em}{0.3em}
\titlespacing*{\subsection}{0pt}{0.5em}{0.2em}

\title{\vspace{-1.2em}\textbf{Visualization System and UI Design for RHFI Framework}\vspace{-0.8em}}
\author{Team 265: Sm Emdadul Hoque}
\date{}

% -----------------------------
\begin{document}
    \twocolumn[
        \maketitle
        \vspace{-1.5em}
    ]

    \section{Visualization Components}

    The visualization system consists of four primary analytical modules designed to support spatial, temporal, comparative, and structural interpretation of housing-market fragility. Each component is directly connected to the RHFI framework and enables multi-dimensional exploration of housing dynamics.

% =========================================================
    \subsection{RHFI Geographic Heatmap}

    The RHFI Geographic Heatmap provides a spatial representation of housing fragility across U.S. states, enabling users to identify regional disparities in housing risk.



    \textbf{Key Features:}
    \begin{itemize}[leftmargin=1.2em]
        \item \textbf{Year-Based Temporal Control:} A slider allows users to dynamically explore changes in housing fragility over time (e.g., pre- and post-2008 financial crisis, COVID-19 period).

        \item \textbf{Multi-Indicator Switching:} Users can toggle between RHFI and underlying indicators such as price growth, rent growth, and inventory change to understand the drivers of regional variation.

        \item \textbf{Quantile-Based Color Scaling:} Instead of fixed ranges, the heatmap applies 5th–95th percentile normalization to reduce the effect of outliers and improve visual stability across years.

        \item \textbf{Rich Hover Tooltips:} Each state displays detailed metrics including RHFI score, affordability ratios, and growth indicators, enabling micro-level interpretation.
    \end{itemize}

    \textbf{Analytical Value:}
    This module enables spatial detection of housing vulnerability clusters, highlighting regions with persistent affordability stress or rapid price acceleration. It is particularly useful for identifying geographic concentration of housing risk.

% =========================================================
    \subsection{Time-Series Explorer}

    The Time-Series Explorer enables longitudinal analysis of housing indicators at the regional level. It is designed to capture structural trends and cyclical behavior in housing markets.

    \textbf{Key Features:}
    \begin{itemize}[leftmargin=1.2em]
        \item \textbf{Long-Term Trend Visualization:} Displays multi-year trajectories of selected indicators such as home prices, rent levels, and RHFI scores.

        \item \textbf{Moving Average Smoothing:} A configurable smoothing window (e.g., 3–5 years) reduces noise and highlights underlying trends.

        \item \textbf{Year-over-Year Change Analysis:} Computes percentage changes to detect acceleration or deceleration in housing pressure.

        \item \textbf{Multi-Metric Comparison:} Allows simultaneous visualization of multiple indicators for structural comparison.
    \end{itemize}

    \textbf{Analytical Value:}
    This module helps identify turning points in housing cycles, including rapid growth phases, stagnation periods, and potential market corrections. It also supports early detection of systemic stress signals.

% =========================================================
    \subsection{Regional Comparison Tool}

    The Regional Comparison Tool enables direct side-by-side comparison of two selected regions, providing a relative assessment of housing fragility.

    \textbf{Key Features:}
    \begin{itemize}[leftmargin=1.2em]
        \item \textbf{RHFI Trajectory Comparison:} Visualizes how housing fragility evolves differently across regions over time.

        \item \textbf{Affordability Gap Analysis:} Compares price-to-income and price-to-rent ratios to evaluate structural affordability differences.

        \item \textbf{Multi-Metric Divergence View:} Displays multiple indicators simultaneously to identify divergence patterns between regions.

        \item \textbf{Interactive Region Selection:} Users can dynamically change regions to explore different geographic pairs (e.g., high-cost vs. emerging markets).
    \end{itemize}

    \textbf{Analytical Value:}
    This component is particularly effective for comparative urban analysis, enabling identification of which regions are becoming more fragile relative to others under similar economic conditions.

% =========================================================
    \subsection{Indicator Breakdown Panel}

    The Indicator Breakdown Panel decomposes the composite RHFI index into its constituent components to improve interpretability and transparency.

    \textbf{Key Features:}
    \begin{itemize}[leftmargin=1.2em]
        \item \textbf{Price Growth Contribution:} Measures contribution of housing price acceleration to overall fragility.

        \item \textbf{Rent Growth Contribution:} Captures rental market pressure and affordability constraints.

        \item \textbf{Inventory Change Impact:} Reflects supply-side dynamics and market tightness.

        \item \textbf{Price-to-Income Ratio Effect:} Represents long-term affordability stress and structural imbalance.
    \end{itemize}

    \textbf{Analytical Value:}
    This decomposition enhances explainability by showing how each macroeconomic factor contributes to the final RHFI score, making the model transparent and interpretable for both technical and non-technical users.

% =========================================================
    \section{User Interface Design}

    The user interface is designed to ensure clarity, accessibility, and analytical efficiency. It supports both exploratory and comparative analysis workflows.

    \textbf{Key UI Design Elements:}
    \begin{itemize}[leftmargin=1.2em]
        \item \textbf{Sidebar Control Panel:} Centralized filtering system for region selection, metric selection, and time control.

        \item \textbf{Tab-Based Navigation:} Structured separation of analytical views (heatmap, time series, comparison, indicators).

        \item \textbf{Interactive Visualization Layer:} All charts support hover interaction, zooming, and dynamic updates.

        \item \textbf{Export Functionality:} Users can download filtered datasets for external analysis.
    \end{itemize}

    \textbf{Design Philosophy:}
    The interface prioritizes minimal cognitive load, ensuring that users can move from high-level insights to detailed analysis without navigating complex menus.

% =========================================================
    \section{Implementation Details}

    The system is implemented using a modern Python-based visualization stack optimized for interactivity and performance.

    \textbf{Core Technologies:}
    \begin{itemize}[leftmargin=1.2em]
        \item \textbf{Streamlit:} Provides the web-based UI framework for rapid dashboard deployment.
        \item \textbf{Plotly Express:} Used for high-level statistical and geographic visualizations.
        \item \textbf{Plotly Graph Objects:} Enables advanced multi-axis and custom interactive charts.
        \item \textbf{Pandas:} Handles data preprocessing, aggregation, and transformation.
        \item \textbf{NumPy:} Supports numerical operations such as normalization and rolling computations.
    \end{itemize}

    \textbf{Performance Optimization:}
    Streamlit caching (\texttt{@st.cache\_data}) is used to reduce redundant computation and improve dashboard responsiveness, especially for large multi-year datasets.

% =========================================================
    \section{System Contributions}

    The visualization system provides three major contributions to the RHFI project:

    \begin{itemize}[leftmargin=1.2em]
        \item \textbf{Transformation of RHFI into an Interactive Tool:} Converts a static index into a dynamic analytical platform.

        \item \textbf{Multi-Dimensional Data Exploration:} Enables spatial, temporal, and comparative analysis within a unified interface.

        \item \textbf{Improved Interpretability and Accessibility:} Makes complex housing-market analytics understandable to non-technical users, including policymakers and researchers.
    \end{itemize}

    \textbf{Overall Impact:}
    The system significantly enhances the usability of the RHFI framework by turning abstract numerical indices into actionable insights through visualization.

% =========================================================


\end{document}